% Собственные макросы и обозначения.
% Класс mipt-thesis-ms уже подгружает amsmath/amssymb, siunitx, esint,
% graphicx, subcaption и др. — здесь только то, чего не хватает.

\usepackage{booktabs} % \toprule/\midrule/\bottomrule для таблиц
\usepackage{tikz}     % схемы (раздел обзора: ось «качество vs стоимость»)
\usetikzlibrary{positioning, arrows.meta, fit, backgrounds}

% --- Типографика и компоновка (CONVENTIONS.md §1: правки только здесь,
%     не в самом .cls) ---
% протрузия улучшает переносы; expansion отключён -- шрифты шаблона (CM)
% не масштабируемые и не поддерживают font expansion
\usepackage[protrusion=true, expansion=false]{microtype}
% Убираем принудительные пустые страницы шаблона: \addemptypage (после
% титула и содержания) и openright-вставки book перед главами. Документ
% читается одностранично, поэтому пустые страницы не нужны.
\renewcommand{\cleardoublepage}{\clearpage}
\renewcommand{\addemptypage}{\clearpage}

% --- Множества чисел ---
\newcommand{\R}{\mathbb{R}}
\newcommand{\N}{\mathbb{N}}
\newcommand{\Z}{\mathbb{Z}}

% --- Векторы, матрицы, множества ---
\newcommand{\vect}[1]{\mathbf{#1}}      % вектор:  \vect{x}
\newcommand{\mat}[1]{\mathbf{#1}}       % матрица: \mat{A}
\newcommand{\set}[1]{\mathcal{#1}}      % множество: \set{X}

% --- Операторы ---
\DeclareMathOperator*{\argmin}{arg\,min}
\DeclareMathOperator*{\argmax}{arg\,max}
\newcommand{\E}{\mathbb{E}}             % мат. ожидание
\newcommand{\norm}[1]{\left\lVert#1\right\rVert}
\newcommand{\abs}[1]{\left\lvert#1\right\rvert}

% Сюда добавляем новые обозначения по мере появления в тексте,
% чтобы система обозначений оставалась сквозной (CONVENTIONS.md §4).

% --- Титульный лист: поля, которых нет в исходном шаблоне ---
% Шаблон mipt-thesis-ms (физтех-школа ФПФЭ) задаёт факультет/кафедру/
% специальность жёстко в \titlepage. Здесь определяем \faculty и
% \department по аналогии с \groupnum и переопределяем \titlepage
% под реквизиты ФПМИ, без правок самого .cls (CONVENTIONS.md §1).
\newcommand{\faculty}[1]{\def\mtfacultytext{#1}}
\newcommand{\department}[1]{\def\mtdepartmenttext{#1}}

\renewcommand{\titlepage}{%
    \pagestyle{empty}
    \begin{center}
        {\bf Министерство образования и науки Российской Федерации \\
        Московский физико-технический институт (национальный
        исследовательский университет) \\ }
        \vspace{1cm}

        \mtfacultytext{} \\
        \mtdepartmenttext{} \\

        \vspace{3em}

        Выпускная квалификационная работа магистра по направлению 09.04.01
        <<Прикладная математика и информатика>>
    \end{center}

    \begin{center}
        \vspace{\fill}
        \singlespacing

        \LARGE \mttitletext{}

        \vspace{\fill}
    \end{center}

    \begin{flushright}
        Студент \mtgroupnumtext{} группы \\
        \mtauthortext{} \\
        \vspace{3em}
        Научный руководитель \\
        \mtsupervisortext{} \\
        \vspace{3em}
        Рецензент \\
        \mtrefereetext{} \\
    \end{flushright}

    \vspace{6em}

    \begin{center}
        Долгопрудный \\
        \the\year{}
    \end{center}
}
